
\section*{Abstract}

We present three Phase 73 components that address systematic biases in CEREBRUM's autonomous knowledge graph discovery pipeline. \textbf{DiscoveryCalibrator} tracks per-community scan and discovery rates via Exponential Moving Average (EMA) and applies an inverse-rate sampling multiplier to steer \texttt{ResearchAgent} toward underrepresented graph regions. \textbf{ContradictionResolver} classifies candidate findings into four evidence classes (clean / revision_candidate / contested / discardable) using a deterministic Noisy-OR evidence model, filtering the most conflicted proposals before they reach \texttt{AutoApprover}. \textbf{CandidateRegistry} replaces the flat evaluated-pairs set with a TTL-aware registry that applies nomination-count boosts and enforces memory bounds via LRU eviction. Together, these components reduce redundant computation, prevent community saturation, and improve the precision of the discovery pipeline by pre-filtering structurally contradicted proposals.

\medskip\noindent\rule{\linewidth}{0.4pt}\medskip

\section*{1. Motivation: Systematic Bias in Autonomous Discovery}

An uncalibrated \texttt{ResearchAgent} exhibits two systematic failure modes:

\begin{enumerate}
  \item \textbf{Community saturation}: The agent repeatedly scans the same densely-connected communities (high community weight) because they yield frequent high-confidence candidates. Underrepresented sparse communities receive no coverage, leaving genuine missing links undiscovered.
\end{enumerate}


\begin{enumerate}
  \item \textbf{Contradiction blindness}: Proposals with strong contradicting evidence in the existing graph (high \texttt{contradiction\_score}) are forwarded to \texttt{AutoApprover} and consume classifier compute, even when deterministic evidence-weight analysis would immediately classify them as discardable.
\end{enumerate}


\medskip\noindent\rule{\linewidth}{0.4pt}\medskip

\section*{2. DiscoveryCalibrator}

\subsection*{2.1 Per-Community EMA Tracking}

For each community \texttt{c}, the calibrator maintains:
\begin{itemize}
  \item \texttt{scan\_rate(c)}: EMA of scans per unit time
  \item \texttt{discovery\_rate(c)}: EMA of approved discoveries per scan
\end{itemize}


\begin{lstlisting}[language=Python]
calibrator.record_scan(community_id)
calibrator.record_discovery(community_id)
\end{lstlisting}

EMA update: \texttt{rate\_t = $\alpha$ $\times$ event + (1 - $\alpha$) $\times$ rate\_\{t-1\}} where \texttt{$\alpha$ = 0.1} (default).

\subsection*{2.2 Inverse-Rate Sampling Multiplier}

The community weight \texttt{w(c)} used in \texttt{\_score\_discovery\_potential()}:

\begin{lstlisting}
w(c) = global_discovery_rate / (discovery_rate(c) + ε)
\end{lstlisting}

Cold-start communities (never scanned) receive \texttt{max\_weight = 5.0}. Communities with higher-than-global discovery rates receive \texttt{w < 1.0} (suppressed). Communities with lower rates receive \texttt{w > 1.0} (boosted).

\subsection*{2.3 Temporal Recency Scoring}

\texttt{ValidationReport.recency\_score} is computed via exponential decay from the publication year:

\begin{lstlisting}
recency_score = exp(-λ × max(0, current_year - pub_year))
\end{lstlisting}

Default \texttt{$\lambda$} corresponds to a 7-year half-life. Recent literature (< 2 years old) scores $\geq$ 0.9.

\medskip\noindent\rule{\linewidth}{0.4pt}\medskip

\section*{3. ContradictionResolver}

\subsection*{3.1 Evidence Model}

For a given finding with \texttt{proposed\_confidences = [c\_1, ..., c\_k]} and \texttt{contradiction\_score}:

\begin{lstlisting}
net_evidence_score = Noisy-OR(c_1, ..., c_k) - contradiction_score
\end{lstlisting}

Noisy-OR: \texttt{1 - ∏(1 - c\_i)} --- the probability of at least one path being correct.

\subsection*{3.2 Resolution Classes}

\begin{table}[h]
\small\centering
\begin{tabular}{lll}
\toprule
net_evidence_score & Class & Action \\
\midrule
$\geq$ 0.6 & \texttt{"clean"} & Forward to AutoApprover \\
0.3 -- 0.6 & \texttt{"revision\_candidate"} & Queue in \texttt{\_revision\_candidates} \\
0.1 -- 0.3 & \texttt{"contested"} & Forward to AutoApprover with penalty \\
< 0.1 & \texttt{"discardable"} & Auto-reject; never reaches AutoApprover \\
\bottomrule
\end{tabular}
\end{table}


Revision candidates accumulate in \texttt{ResearchAgent.\_revision\_candidates} for periodic batch review.

\medskip\noindent\rule{\linewidth}{0.4pt}\medskip

\section*{4. CandidateRegistry}

\subsection*{4.1 TTL-Aware Registry}

Replaces the flat \texttt{\_evaluated\_pairs: Set[Tuple[str, str]]} with a dict-based registry:

\begin{lstlisting}[language=Python]
registry[pair] = CandidateRecord(
    nomination_count=N,
    first_seen=t_0,
    last_seen=t_N,
    ttl=timedelta(hours=24),
)
\end{lstlisting}

\texttt{is\_registered(pair)} returns True and blocks re-evaluation if within TTL.

\subsection*{4.2 Nomination Boost}

When a pair is nominated multiple times (by different scan cycles or reasoning strategies), its \texttt{discovery\_potential} receives a log-scale boost:

\begin{lstlisting}
boosted_potential = raw_potential × min(log(N + 1) + 1, max_boost)
\end{lstlisting}

Default \texttt{max\_boost = 3.0}. This rewards repeatedly-surfaced candidates without linearly amplifying noise.

\subsection*{4.3 Memory Bound}

\texttt{max\_entries} (default 10,000) triggers LRU eviction: the registry evicts the least-recently-seen entry when capacity is reached, ensuring bounded memory usage in long-running deployments.

\medskip\noindent\rule{\linewidth}{0.4pt}\medskip

\section*{5. References}

\begin{itemize}
  \item [pearl2000causality] Pearl, J. (2000). \textit{Causality: Models, reasoning, and inference.} Cambridge University Press.
\end{itemize}
